\chapter{The Servants}

\section{The Problem With A General Referent}

The agency chapter changed the servant question. If the familiar chain is assumed, the servants are easy to identify in broad terms. God gives the disclosure to Jesus Christ; Jesus Christ shows it to his servants; John receives the message; the churches hear it. On that reading, ``his servants'' naturally becomes a general description of Christians who receive the book.

That reading is understandable. The book of Revelation is addressed to churches. The blessing in Revelation 1:3 includes the one who reads, the ones who hear, and the ones who keep what is written. The later book repeatedly calls God's faithful people servants. A broad Christian referent is therefore not invented from nothing.

The problem is the word ``generally.'' The phrase in Revelation 1:1 is \textgreek{τοῖς δούλοις αὐτοῦ} (\textit{tois doulois autou}, \textit{to his servants}). It does not simply say ``to Christians.'' It does not yet name the churches. It does not yet introduce the reader and hearers of verse 3. It names servants within the opening communication sequence.

That difference matters because \textit{doulos} is a relational term. A servant is not defined only by belonging to a religious group. A servant stands in relation to a master, lord, or ruler. The term carries a frame of command, reception, loyalty, and obedience. If the opening says that God shows and signifies the things that must happen to his servants, then the servants are not merely an audience who learns information. They are subjects addressed within a relation of service.

The possessive pronoun also matters. The phrase is not bare: ``the servants.'' It is ``his servants.'' The preceding chapter argued that God is the giver, shower, and signifier. On that analysis, the possessive most naturally belongs to God. They are God's servants, not because later Christian theology says believers belong to God in a general sense, but because the local communication structure places God as the acting master who gives, shows, signifies, and sends.

The same opening also calls John \textgreek{τῷ δούλῳ αὐτοῦ Ἰωάννῃ} (\textit{to doulo autou Ioanne}, \textit{to his servant John}). This is not a minor repetition. The plural servants in the C phrase and the singular servant John in the B' phrase belong to the same lexical field. The plural are the recipients of what is shown and signified. John is the named servant to whom the disclosure is sent through the angel. Any interpretation of the plural servants must explain why the same opening names John with the singular form of the same term.

This prevents two shortcuts. The first shortcut is to make the servants simply all Christians as a flat category. The second is to isolate John from the servant category, as though he were only a neutral courier standing outside the relation named by the text. The opening does neither. It places John and the plural servants within one master-servant frame, while still distinguishing their roles in the communication sequence.

The point is not that later Christian readers are excluded. That would overstate the evidence. The point is that later readers must enter the passage through the role the passage names. The servants are not first defined by later reception of the book. They are first defined by the opening's own structure: God's servants are the ones to whom the things that must happen are shown and signified, while John is the servant through whom the disclosure becomes testimony.

\section{Servants In The Book Of Revelation}

The wider book confirms that servant language is not a casual synonym for religious audience. It is flexible, but it is not empty.

At some points, servants are God's protected people. Revelation 7:3 speaks of the servants of God being sealed. At other points, the language is explicitly prophetic. Revelation 10:7 refers to God's servants the prophets. Revelation 11:18 brings together ``your servants the prophets,'' the saints, and those who fear God's name. At the end of the book, the angel can call himself a fellow servant with John and with John's brothers who hold the testimony or keep the words of the book. Revelation 22:6 then returns to language close to the opening: God sent his angel to show his servants what must happen.

Those later uses do not reduce the phrase in Revelation 1:1 to one narrow class. They do show the kind of field in which the word operates. Servants are those who belong to God, receive what God communicates, bear witness, keep the words, worship, suffer, and stand under divine command. The term can include prophets, saints, hearers, keepers, and worshipers, but it includes them as servants, not as an undefined crowd.

This wider pattern makes three possible referents worth testing.

First, the servants may be prophetic servants. That option has real strength. John is called a servant. John testifies. The book calls prophets God's servants. The opening sequence is a communication sequence, and prophecy is the form in which the written book is later read and kept. If the servant category is controlled by testimony and prophetic reception, then the servants are not simply every person who happens to read the book. They are those who stand in the prophetic service of receiving, preserving, and obeying God's communicated word.

But that option becomes too narrow if it excludes the churches and hearers who are explicitly blessed in verse 3. The book of Revelation does not reserve obedience and witness for a technical prophetic class. The one reading, the ones hearing, and the ones keeping are drawn into the same communicative event. The prophetic frame governs the mode of reception, but it need not restrict the servants only to official prophets.

Second, the servants may be the churches addressed by the book. This also has strength. Revelation 1:4 soon addresses the seven churches, and the messages in chapters 2--3 show that the written work is not detached from real congregations. The book is meant to be read, heard, and kept among them. If the servants are God's obedient people in those churches, then the phrase reaches the public audience of the book.

But this option becomes too broad if it is stated as ``Christians generally'' without qualification. The churches are not addressed as a religious demographic. They are summoned to hear, keep, repent, endure, conquer, and remain faithful. In other words, they are addressed under the demands of service. The servant category may include the churches, but it interprets their identity rather than merely naming their membership.

Third, the servants may include heavenly or angelic servants. The book's angel can call himself a fellow servant, so the category is not limited in every context to human believers. Yet Revelation 1:1 immediately moves toward John, testimony, reading, hearing, and keeping what is written. The controlling line of reception in the opening is therefore human and prophetic, not primarily angelic. Angelic service belongs to the book's larger world, but it does not supply the best immediate referent for \textit{tois doulois autou} in the opening sentence.

The most disciplined conclusion is therefore not a single flat label. The servants in Revelation 1:1 are God's obedient recipients within the disclosure system. John is the named servant who receives the disclosure through the angel and testifies. The plural servants are those to whom God shows and signifies the things that must happen, and who are thereby summoned into faithful reception and obedience. The later churches and readers can be included only as they occupy that servant role.

\section{The Servants In The Communication Chain}

The servant question must finally be placed back into the expanded form. The C / C' pair supplies the same recipients and content to both showing and signifying:

\begin{quote}
\textit{deixai tois doulois autou ha dei genesthai en tachei}
\end{quote}

\begin{quote}
\textit{kai [tois doulois autou ha dei genesthai en tachei] esemanen}
\end{quote}

The servants are therefore not incidental. They are the audience of two paired communication acts. The things that must happen are shown to them, and on this analysis the same things are signified to them. The servant identity determines the kind of reception the passage imagines. These recipients do not merely observe information about the future. They receive a disclosure from their master in a form that calls for recognition, trust, and obedience.

John's role is related but not identical. The B / B' pair identifies John as the recipient of what God gives and sends through the angel. John then testifies to the word of God, the testimony of Jesus Christ, and what he saw. He is not outside the servant category; he is explicitly called God's servant. But he is the named servant-witness through whom the disclosure becomes testimony for others. The plural servants are the intended obedient recipients of what God shows and signifies.

This means the opening does not give a simple chain in which the servants sit passively at the end. The structure is more layered. God gives the disclosure to John. God shows and signifies the things that must happen to his servants. God sends the disclosure through his angel to his servant John. John testifies. The reader and hearers are then blessed if they keep what is written. The servants are embedded in this whole movement, not merely appended to it.

That layered structure explains why ``all Christians generally'' is insufficient. If the phrase only meant a broad later audience, it would not explain the relation between the plural servants and John the servant. It would not explain why showing and signifying are directed to servants before verse 3 introduces reading, hearing, and keeping. It would not explain why the term carries a master-servant frame at the exact point where God is acting as giver, shower, signifier, and sender.

A better formulation is this: the servants are God's loyal subjects addressed by the disclosure, with John as the named servant-witness and the churches/readers included insofar as they receive and keep the prophetic word as servants. This keeps the category broad enough to reach the book's public audience, but controlled enough to preserve the opening structure.

The consequence is important. What God shows and signifies is not spectacle for curious observers. It is disclosure for servants. That means the phrase that follows, \textit{ha dei genesthai en tachei}, must be read in relation to servant reception. The next question is not merely whether the events are soon. It is what kind of necessity and speed belong to things shown and signified to God's servants within this command-and-obedience frame.
